% ============================================================
% BME VIK TDK abstract — VARIANT 2: explainability / capability-
% boundary framing ("does the LLM actually use the injected graph
% signal, and what does it add over text alone?")
%
% This is the second of two abstract framings for the supervisors to
% compare and choose between. Variant 1 (Writing/abstract-v1-problem-
% method.tex) frames the work as reconstructing a paper's abstract
% from its citation graph (task-first). This variant instead treats
% that same generation task as a testbed for a behavioural/diagnostic
% question about graph-conditioned LLMs in general.
%
% This framing operationalises the already-locked primary TDK
% narrative -- per tdk-pitch-reframe.md's recommended formal title
% and tdk-winning-strategy.md's "one-line strategy" (lead with the
% falsifiable question: does the graph signal survive, or is the
% model just reading the text?).
%
% Structure follows the official one-page sample:
% https://portal.vik.bme.hu/files/00000047.pdf
%   - max 1 page, min ~20 lines of body text
%   - NO header, footer, or page numbers
%   - figures/equations only if strictly necessary
%   - own contribution must be clearly separated from cited work
%   - body must cover: methodology, key results, conclusions
%
% NOTE ON LANGUAGE: the official TDK conference language is Hungarian
% (English talks require prior agreement with the section chair, per
% the official sample). This English draft is for internal review with
% Bálint/Dani; translate to Hungarian before submission unless an
% English exception is confirmed.
% ============================================================

\documentclass[12pt]{article}

\usepackage[utf8]{inputenc}
\usepackage[T1]{fontenc}
\usepackage[english]{babel}
\usepackage[a4paper, margin=2.5cm]{geometry}
\usepackage{mathptmx}      % Times-like font
\usepackage{hyperref}
\usepackage{enumitem}

\pagestyle{empty}          % <- enforces "no header/footer/page numbers"
\setlength{\parindent}{0pt}
\setlength{\parskip}{0.6em}

\begin{document}

% ---------- Title ----------
\begin{center}
    {\Large\bfseries Is the Language Model Actually Listening to the Graph? \\
    Structure-Conditioned Abstract Generation on Citation Networks}
\end{center}

\vspace{0.8em}

% ---------- Authors ----------
\begin{center}
    Gema Csanád IV. Inf, gemacsanad@gmail.com \\

\end{center}

\vspace{0.4em}

% ---------- Consultant ----------
\begin{center}
    Consultant: Unyi Dániel
\end{center}

\vspace{1.2em}

% ==========================================================
% BODY
% ==========================================================

A citation graph represents scientific literature as a network: each
node is a publication, each edge is a citation between two papers,
and every node carries that paper's abstract as its text content. One
way to condition a language model on this graph is to encode a
paper's citation neighbourhood into a continuous embedding and inject
it into the model as an additional signal, alongside the neighbours'
abstracts, which are already given to the model as plain text. This
raises a question existing work does not answer: does the language
model actually use this injected signal, or does it default to
relying on the text that is already there? We study this question
directly, using the task of generating a paper's abstract from its
citation neighbourhood as the testbed.

We use a graph neural network (GNN) to encode the citation
neighbourhood into a structural embedding, project it through a
trained adapter into soft prefix tokens, and inject these into every
attention layer of a LoRA-fine-tuned language model, alongside the
plain-text neighbour abstracts. A contrastive alignment loss between
the GNN embedding and a frozen text encoding of the true target
abstract is added during training, both to encourage the embedding to
carry useful information and to serve as an internal check on whether
that information is present at all before it ever reaches the
language model. Isolating the injected signal's contribution requires
more than measuring final generation quality: we compare the full
model against a text-only baseline that receives no injected signal
at all, and against a variant that receives the same soft prefix but
derived from a shuffled, incorrect neighbourhood, to test whether the
model responds to the prefix's specific content or merely to its
presence.

We construct the task on ogbn-arxiv-TA, a citation graph of computer
science papers extended with cross-domain reference nodes, and
evaluate generation quality with ROUGE, rescaled BERTScore, and an
LLM-as-judge protocol with explicit controls for position, verbosity,
and self-preference bias.

Prior graph-conditioned language models report only end-task
performance, treating the graph's contribution as a black box; none
test whether an injected structural signal is actually used before
crediting it with any improvement. Our ablations are designed to
answer this directly, and we treat both possible outcomes as
informative results, not as success or failure of the project: if the
graph-conditioned model diverges from the shuffled-neighbour and
text-only variants, this is evidence that the language model can
extract and use structural information beyond what the neighbours'
own text already conveys; if it does not, this is evidence of a
text-dominance effect, in which the model defaults to the readily
available text and never learns to rely on the injected channel -- a
finding with implications for how graph-conditioned language models
should be evaluated more broadly, beyond this one task.

\end{document}
